\chapter{Witt theory of quadratic spaces}
\label{ch:la01-witt}

Throughout this chapter \(k\) is a field with
\(\operatorname{char} k\ne 2\), and the ambient quadratic spaces are
finite dimensional and nondegenerate. For a quadratic form \(q\) we
use the symmetric bilinear form
\[
  B(x,y)=\frac{q(x+y)-q(x)-q(y)}{2},
  \qquad q(x)=B(x,x).
\]
Nondegeneracy means that \(B\) has zero radical. An isometry is a
linear isomorphism preserving \(q\), equivalently \(B\). We write
\(\perp\) for an orthogonal direct sum and \(\langle a\rangle\) for
the line with form \(x\mapsto ax^2\). These conventions agree with
the polarization and diagonalization of the archived notes
\cite[11.polarization and 11.diagonal\_congruence]{d05}.

For a subspace \(U\subseteq V\), write
\(U^\perp=\{x\in V:B(x,U)=0\}\).
When the restriction to \(U\) is nondegenerate, the decomposition
\(V=U\perp U^\perp\) is the symmetric case of the projection argument
in \cite[11.nondegenerate]{d05}. In fact, for \(v\in V\), the
restriction \(B(v,-)|_U\) has a unique representative \(u\in U\)
under \(U\cong U^*\) via \(B|_U\). Thus \(v-u\in U^\perp\).
The restriction to \(U^\perp\) is nondegenerate because a vector
orthogonal to both summands belongs to the radical of \(V\).

\section{Isotropic vectors and hyperbolic planes}

A nonzero vector \(e\) is \emph{isotropic} if \(q(e)=0\).
A space is \emph{anisotropic} when its only vector of value zero is
the zero vector, so the zero space is anisotropic. A subspace \(T\)
is \emph{totally isotropic} if \(q|_T=0\); by polarization this is
equivalent to \(B|_{T\times T}=0\). The hyperbolic plane \(H\) is
\(k^2\) with \(q_H(x,y)=xy\), equivalently
\(\langle 1\rangle\perp\langle -1\rangle\) by the invertible
change \((s,t)\mapsto(s+t,s-t)\).

\begin{lemma}[LA01-C06.01. Splitting a hyperbolic plane]
\label{lem:la01-hyperbolic-splitting}
If \(e\in V\) is isotropic, there is an isotropic vector \(f\)
with \(2B(e,f)=1\). Then \(P=ke+kf\) is isometric to \(H\) and
\(V=P\perp P^\perp\) with \(P^\perp\) nondegenerate.
\end{lemma}

\begin{proof}
Nondegeneracy and \(e\ne 0\) give a vector pairing nontrivially
with \(e\). Scale it to \(f_0\) with \(2B(e,f_0)=1\), and put
\(f=f_0-q(f_0)e\). The identities
\[
  2B(e,f)=1,\qquad
  q(f)=q(f_0)-2q(f_0)B(e,f_0)+q(f_0)^2q(e)=0
\]
show that \(e,f\) are independent and \(q(se+tf)=st\).
Their Gram matrix has diagonal zero and off-diagonal entries
\(1/2\), so the restriction to \(P\) is nondegenerate. Apply the
orthogonal splitting for a nondegenerate restriction.
\end{proof}

\section{Reflections and cancellation}

\begin{lemma}[LA01-C06.02. Moving vectors of the same nonzero value]
\label{lem:la01-reflections}
For \(w\in V\) with \(q(w)\ne 0\), the reflection
\[
  s_w(x)=x-\frac{2B(x,w)}{q(w)}\,w
\]
is an isometry. If \(u,v\in V\) satisfy
\(q(u)=q(v)=a\ne 0\), an isometry sending \(u\) to \(v\)
can be chosen as a product of at most two reflections.
\end{lemma}

\begin{proof}
The orthogonal decomposition \(V=kw\perp(kw)^\perp\) shows
that \(s_w\) negates the \(kw\) coordinate and fixes
\((kw)^\perp\). It is therefore an involution and preserves \(q\).
If \(q(u-v)\ne 0\), then \(2B(u,u-v)=q(u-v)\), so
\(s_{u-v}(u)=v\). Otherwise
\[
  q(u+v)=2q(u)+2q(v)-q(u-v)=4a\ne 0.
\]
In this case \(2B(u,u+v)=q(u+v)\), so \(s_{u+v}(u)=-v\).
The reflection \(s_v\) sends \(-v\) to \(v\). This second case
also works when \(u=v\).
\end{proof}

\begin{theorem}[LA01-C06.03. Witt cancellation]
\label{thm:la01-witt-cancellation}
Let \(U,V,W\) be nondegenerate quadratic spaces. If
\(U\perp W\cong V\perp W\), then \(U\cong V\).
The same conclusion holds when the copies of \(W\) are
replaced by two isometric nondegenerate spaces.
\end{theorem}

\begin{proof}
First let \(W=ke\) with \(q(e)=a\ne 0\). For an isometry
\(F:U\perp W\to V\perp W\), the vectors \(F(e)\) and the
specified \(e\) in the target both have value \(a\).
By Lemma~\ref{lem:la01-reflections} there is a target isometry
\(s\) sending \(F(e)\) to the target \(e\). The isometry \(sF\)
then carries the orthogonal complement of the source line to
the orthogonal complement of the target line. These
complements are \(U\) and \(V\).

For arbitrary \(W\), take an orthogonal basis by
\cite[11.diagonal\_congruence]{d05}. Nondegeneracy makes every
diagonal coefficient nonzero. Cancel the lines one at a time
by the preceding argument. When \(W=0\), the given isometry is
already the desired isometry. For two isometric copies of
\(W\), identify them by their given isometry before cancelling.
\end{proof}

\section{Witt decomposition and Witt index}

\begin{theorem}[LA01-C06.04. Witt decomposition and uniqueness]
\label{thm:la01-witt-decomposition}
There is an isometry
\[
  V\cong H^{\perp r}\perp A
\]
for an integer \(r\ge 0\) and an anisotropic space \(A\). The
integer \(r\) and the isometry class of \(A\) depend only on \(V\).
We call \(A\) the anisotropic part and \(r\) the Witt index.
Here \(H^{\perp 0}=0\).
\end{theorem}

\begin{proof}
If \(V\) is anisotropic, take \(r=0\). Otherwise split off a
hyperbolic plane by Lemma~\ref{lem:la01-hyperbolic-splitting} and
repeat in its nondegenerate complement. The dimension drops by
two at each step, so this process ends with an anisotropic
remainder, possibly zero.

For another decomposition \(H^{\perp s}\perp A'\), assume
\(r\le s\) after exchanging the two decompositions if needed.
Witt cancellation (Theorem~\ref{thm:la01-witt-cancellation})
gives \(A\cong H^{\perp(s-r)}\perp A'\). If \(s-r>0\), the
right side has a nonzero isotropic vector, contradicting
anisotropy of \(A\). Thus \(r=s\) and \(A\cong A'\).
\end{proof}

\begin{proposition}[LA01-C06.05. Totally isotropic subspaces and the Witt index]
\label{prop:la01-witt-index}
Let \(T\subseteq V\) be totally isotropic of dimension \(m\).
There is an orthogonal sum of \(m\) hyperbolic planes in \(V\)
whose chosen first isotropic basis vectors span \(T\).
Consequently \(m\le r\), and \(r\) is the maximum dimension of a
totally isotropic subspace. Every totally isotropic subspace
maximal by inclusion has dimension \(r\).
\end{proposition}

\begin{proof}
For \(m=0\) take the empty sum. For \(m>0\), choose
\(e\in T\setminus\{0\}\) and use
Lemma~\ref{lem:la01-hyperbolic-splitting} to write
\(V=P\perp P^\perp\) with \(P=ke+kf\) and \(2B(e,f)=1\).
For \(t\in T\), total isotropy gives \(B(t,e)=0\). The vector
\[
  t-2B(t,f)e
\]
also pairs to zero with \(f\), so it lies in \(T\cap P^\perp\).
Thus \(T=ke\oplus(T\cap P^\perp)\). Apply the induction
hypothesis to this \((m-1)\)-dimensional subspace inside the
nondegenerate space \(P^\perp\). This produces the claimed \(m\)
planes, with a nondegenerate orthogonal complement \(C\).

Decompose \(C\) by Theorem~\ref{thm:la01-witt-decomposition}.
Uniqueness gives \(r=m\) plus the Witt index of \(C\), hence
\(m\le r\). The span of one isotropic coordinate vector in each
of the \(r\) planes in a Witt decomposition has dimension \(r\).
If \(T\) is maximal by inclusion, \(C\) is anisotropic, since an
isotropic line in \(C\) could be adjoined orthogonally to \(T\).
Then the Witt index of \(C\) is zero and \(m=r\).
\end{proof}

For the general theory, see \cite[Ch.~II, Secs.~7--8]{witt-source}.

\begin{example}[LA01-C06.06. A Witt index that increases after extending scalars]
\label{ex:la01-witt-base-change}
On \(\mathbb{Q}^4\) take \(q(x,y,u,v)=xy+u^2-2v^2\).
The last two coordinates form a nondegenerate plane with no
nonzero isotropic vector. Indeed, \(u^2=2v^2\) with \(v\ne 0\)
would make \(2\) a rational square: writing \(u/v=a/b\) in
lowest terms would give \(a^2=2b^2\), forcing both \(a\) and
\(b\) even. If \(v=0\), then \(u=0\). The displayed form is
therefore \(H\) plus an anisotropic plane, and its Witt index
over \(\mathbb{Q}\) is \(1\).

Over \(\mathbb{Q}(\sqrt{2})\) the change
\((u,v)\mapsto(u+\sqrt{2}v,u-\sqrt{2}v)\) is invertible and
turns \(u^2-2v^2\) into a product of the two new coordinates.
That plane becomes hyperbolic, so the Witt index after extension
is \(2\). Both computations use
Theorem~\ref{thm:la01-witt-decomposition}.
\end{example}
